\documentclass[11pt]{article}
\usepackage[margin=1in]{geometry}
\usepackage{graphicx}
\usepackage{hyperref}
\usepackage{booktabs}
\usepackage{float}
\usepackage{amsmath}
\usepackage{setspace}

\onehalfspacing

\title{\textbf{Transaction Fraud Classifier: Project Report}}
\author{Nathaniel Wetzel}
\date{April 20, 2026}

\begin{document}

\maketitle

\section{Problem Description}

\subsection{Machine Learning Problem}
For my final project, I decided to address the problem of credit card fraud detection. In machine learning terms, this constitutes a \textbf{binary classification problem}. I analyzed historical credit card transactions to predict whether a new, incoming transaction is legitimate (class 0) or fraudulent (class 1). The model takes in features about the transaction (such as transaction amount, location, and user details) and outputs a probability score indicating the likelihood of fraud.

\subsection{Motivation}
I chose this problem because credit card fraud is a massive real-world issue. According to various financial reports, billions of dollars are lost every year due to unauthorized transactions. When a fraudster steals someone's card information, they generally try to maximize the financial extraction quickly before the card is blocked. 

From a machine learning perspective, this problem is a compelling and complex challenge because the data is extremely imbalanced. Most people are honest, so out of millions of transactions, only a tiny fraction represent actual fraud. I wanted to see how different algorithms handle this extreme class imbalance and determine if it was possible to build a system that detects fraud without accidentally declining legitimate purchases (which is highly disruptive to cardholders).

\subsection{References and Related Work}
When researching how to approach the dataset, I reviewed several key resources and related datasets (such as the popular IEEE-CIS Fraud Detection dataset and the ULB Machine Learning Group dataset) to learn how others have solved this:
\begin{enumerate}
    \item \textbf{A Survey of Credit Card Fraud Detection Techniques (Sorournejad et al., 2016)}: This paper was highly useful for understanding the history and complexities of the problem. It surveyed the state of the art in credit card fraud detection, discussing the difference between supervised (misuse) and unsupervised (anomaly) approaches, as well as the unique challenges of working with extremely imbalanced datasets. This reinforced my decision to treat the problem as a supervised anomaly detection task where generating features that capture behavioral shifts is paramount. \url{https://arxiv.org/abs/1611.06439}
    \item \textbf{Kaggle Community Notebooks on Fraud Detection}: Since I sourced my data from Kaggle, I investigated how other practitioners approached similar tabular datasets. Many top-performing implementations relied heavily on tree-based models like XGBoost and LightGBM. I learned that decision trees are robust to unscaled data, but for my neural network, I would definitely need to normalize my inputs. This led directly to my implementation of a `StandardScaler` pipeline. \url{https://www.kaggle.com/datasets/kartik2112/fraud-detection/code}
    \item \textbf{Deep Learning for Imbalanced Data}: I read a few articles detailing how Neural Networks generally struggle with tabular data compared to image or text data. Specifically, Multi-Layer Perceptrons (MLPs) can easily collapse to majority class predictions (predicting "not fraud" every time) if metrics and architecture are not carefully managed. This informed my central decision to evaluate all models using F1-scores and ROC-AUC rather than relying on baseline accuracy. \url{https://www.tensorflow.org/tutorials/structured_data/imbalanced_data}
\end{enumerate}

\section{Dataset}

\subsection{Dataset Statistics}
I used the Kaggle "Fraud Detection" dataset by user kartik2112 (\url{https://www.kaggle.com/datasets/kartik2112/fraud-detection}), which simulates credit card transactions from a variety of users over a long period. 

Here are the basic statistics of the data I worked with:
\begin{itemize}
    \item \textbf{Number of Records}: When I combine the train and test files, I am working with over 1.8 million transactions. 
    \item \textbf{Classes}: There are exactly two classes in my target variable `is\_fraud`. Legitimate transactions are marked as 0, and fraudulent ones are marked as 1.
    \item \textbf{Number of Features}: The raw dataset comes with 22 columns. After dropping identifiers (like the transaction ID or raw names), I used a core set of about 8 raw features, plus 4 custom features I engineered later on.
    \item \textbf{Feature Types}: The data includes continuous numerical features (amount, city population), categorical text variables (merchant category, user gender, state of residence), and datetime variables.
    \item \textbf{Missing Data}: The dataset was exceptionally clean. Unlike real-world messy data, this simulated dataset didn't have missing values in the core columns, so I didn't have to perform complex imputation like K-Nearest Neighbors. I simply dropped any anomalous rows during the data loading phase.
\end{itemize}

\subsection{Feature Definition and Insights}
The main features I decided to focus on from the raw data were:
\begin{itemize}
    \item \textbf{amt}: The transaction amount. In my EDA, I noticed fraudulent activity often begins with small probing transactions followed by a massive purchase.
    \item \textbf{category}: The category of the merchant (like 'grocery\_pos', 'entertainment', 'shopping\_net'). Fraudsters often prefer online shopping because they don't have to show physical identification.
    \item \textbf{city\_pop}: The population of the city. I included this hypothesis testing whether rural areas have different baseline fraud rates than major cities.
    \item \textbf{Time features}: The raw data included a `trans\_date\_trans\_time` string. I parsed this into `trans\_hour`, `trans\_dayofweek`, and `trans\_month`. Fraud tends to occur at irregular hours, such as 3 AM, when the actual cardholder is likely asleep.
    \item \textbf{Demographics}: The `gender` and `state` of the cardholder. I included these to help the model identify localized trends.
\end{itemize}

\subsection{Exploratory Data Analysis}
Before feeding data into the models, I did some exploratory data analysis (EDA). The very first thing I checked was the class balance, because I knew fraud data is notoriously skewed. 

\begin{figure}[H]
    \centering
    \begin{minipage}{0.48\textwidth}
        \centering
        \includegraphics[width=\linewidth]{outputs/class_balance.png}
        \caption{Class Balance demonstrating extreme rarity of fraud.}
        \label{fig:class_balance}
    \end{minipage}\hfill
    \begin{minipage}{0.48\textwidth}
        \centering
        \includegraphics[width=\linewidth]{outputs/amount_distribution.png}
        \caption{Distribution of transaction amounts by class (truncated to < \$1000).}
        \label{fig:amount_distribution}
    \end{minipage}
\end{figure}

As Figure \ref{fig:class_balance} shows, the number of legitimate transactions completely dwarfs the fraudulent ones. In my code, I calculated that fraud makes up exactly 0.52\% of the entire dataset. This immediately confirmed my suspicion that standard accuracy would be an ineffective metric. 

Additionally, as seen in Figure \ref{fig:amount_distribution}, the distribution of transaction amounts varies significantly between classes. Legitimate transactions are heavily concentrated in smaller amounts (under \$100), whereas fraudulent transactions show higher density in the \$200--\$400 range and exhibit an extended long tail, confirming that `amt` and moving averages of `amt` are highly predictive features.

\section{Approach and Methodology}

\subsection{General Approach}
My general pipeline was to first construct a baseline using just the raw features. I developed an automated pipeline (`fraud\_detection.py`) to systematically preprocess the data, evaluate it across four distinct machine learning algorithms, and report evaluation metrics. Once establishing this baseline, the primary objective was to determine the exact aggregate performance difference achieved by introducing my novel custom features and running the identical pipeline side-by-side.

\subsection{Feature Selection, Normalization, and Defining New Features}
Machine learning models, especially Neural Networks and Logistic Regression, rely on optimization functions that assume standardized feature magnitudes. Without standardization, passing in a transaction amount of \$5,000 and a `trans\_hour` of 2 would lead the model to disproportionately weigh the transaction amount.

To fix this, I used \textbf{StandardScaler} from scikit-learn on all my numeric continuous columns (`amt`, `city\_pop`, etc.). This forces the data to have a mean of 0 and a standard deviation of 1 using the formula:
\begin{equation}
z = \frac{x - \mu}{\sigma}
\end{equation}

For text data like `category`, `gender`, and `state`, models can't read strings. I considered using One-Hot Encoding, but I realized that since there are 50 states and lots of categories, One-Hot Encoding would create way too many columns (the "curse of dimensionality") and slow my training down to a crawl. Instead, I used \textbf{LabelEncoder}, which just assigns an integer to each category (e.g., NY = 1, CA = 2). It's not perfect since states don't have a mathematical order, but tree-based models handle this perfectly fine, and it kept my code fast.

Regarding defining new features, I chose to construct new temporal "velocity" variables (such as counting the number of transactions a card made in the last 24 hours). This is detailed further in the Novel Contributions section, as it forms the core of my unique additions to the dataset.

\subsection{Machine Learning Models}
I chose four different models to get a good mix of simple, advanced, and deep learning algorithms:

\begin{enumerate}
    \item \textbf{Logistic Regression}: This was my baseline model. It operates as a linear model, attempting to draw a linear decision boundary between fraud and non-fraud. I didn't expect it to perform exceptionally well, but I needed a robust baseline to compare against.
    \item \textbf{Random Forest}: This is an ensemble method that builds numerous decision trees and takes a majority vote. Because it partitions data on thresholds, I knew it would be well-suited for handling my Label-Encoded categorical data.
    \item \textbf{XGBoost}: This is an extreme gradient boosting model involving sequential tree building, where each new tree aims to correct the residual errors of the previous ones. Based on the academic literature and Kaggle discussions, this is a leading algorithm for tabular structured data.
    \item \textbf{Neural Network (MLPClassifier)}: I wanted to include a deep learning approach, so I configured a Multi-Layer Perceptron. I used the default hidden layer architectures but relied on my StandardScaler to aid convergence. I aimed to see if a neural net could uncover non-linear interactions better than XGBoost.
\end{enumerate}

Because my dataset is incredibly imbalanced (99\% legitimate), if a model predicts "legitimate" for every transaction, it yields 99\% Accuracy. Therefore, Accuracy is an ineffective metric for this domain. Instead, I evaluated my models using Precision, Recall, F1-Score, and ROC-AUC. 
F1-Score is the harmonic mean of Precision and Recall. This provided a single robust metric to easily compare the models.

\subsubsection{Comparison of the Models}
After completing both pipeline evaluations, the empirical results provided an excellent contrast in algorithmic capabilities. The following figures summarize both the aggregate F1 performance across different feature sets and the underlying receiver operating characteristics.

\begin{figure}[H]
    \centering
    \includegraphics[width=0.8\textwidth]{outputs/comparison_f1.png}
    \caption{Comparison of F1 Scores across models. Notice the massive jump when I add my novel temporal features.}
    \label{fig:f1_compare}
\end{figure}

Looking at Figure \ref{fig:f1_compare}, it is evident that XGBoost and Random Forest significantly outperformed the Logistic Regression and Neural Network models. Logistic Regression barely achieved a 0.128 F1 score on the raw data, indicating it could not map the complexity of the feature space. More importantly, for every single model, adding my custom transaction velocity features generated a substantial improvement in performance. XGBoost improved from an F1 score of 0.674 up to 0.892 just by incorporating the temporal transaction variables. 

\begin{figure}[H]
    \centering
    \includegraphics[width=0.85\textwidth]{outputs/roc_Raw_+_Temporal_Features.png}
    \caption{ROC Curve for all models trained on my enhanced Temporal dataset.}
    \label{fig:roc_temporal}
\end{figure}

The ROC curve (Figure \ref{fig:roc_temporal}) shows that almost all the models (except Logistic Regression) are extremely competent at separating the classes based on probability distributions, achieving over 0.95 AUC. 

\subsubsection{Discussion of Results and Insights}

To practically understand where the models were making errors, I analyzed the confusion matrices from the best performing feature set (Raw + Temporal).

\begin{figure}[H]
    \centering
    \begin{minipage}{0.45\textwidth}
        \centering
        \includegraphics[width=\linewidth]{outputs/cm_Raw_+_Temporal_Features_Logistic_Regression.png}
        \caption{Logistic Regression CM}
    \end{minipage}\hfill
    \begin{minipage}{0.45\textwidth}
        \centering
        \includegraphics[width=\linewidth]{outputs/cm_Raw_+_Temporal_Features_Neural_Network_MLP.png}
        \caption{Neural Network CM}
    \end{minipage}
\end{figure}

The Logistic Regression model (left) predicted a high volume of False Positives and False Negatives. The Neural Network (right) performed much better at catching True Negatives, but it missed actual instances of fraud (higher False Negatives), meaning it was penalizing the minority class and predicting too conservatively. It is crucial to note that the Neural Network achieved an exceptional ROC-AUC of 0.997, technically matching XGBoost. However, it struggled on the F1 score, demonstrating that while the network learned a highly accurate probability distribution (AUC), its default classification threshold ($p=0.5$) was poorly calibrated for this extreme minority class. Tabular deep learning typically requires threshold-moving or extensive hyperparameter tuning to translate high AUCs into high F1 scores.

\begin{figure}[H]
    \centering
    \begin{minipage}{0.45\textwidth}
        \centering
        \includegraphics[width=\linewidth]{outputs/cm_Raw_+_Temporal_Features_Random_Forest.png}
        \caption{Random Forest CM}
    \end{minipage}\hfill
    \begin{minipage}{0.45\textwidth}
        \centering
        \includegraphics[width=\linewidth]{outputs/cm_Raw_+_Temporal_Features_XGBoost.png}
        \caption{XGBoost CM}
    \end{minipage}
\end{figure}

The tree-based models were highly robust. XGBoost (right) caught the vast majority of the actual fraud cases while maintaining incredibly low False Positives. If a real financial institution implemented this XGBoost model with the temporal features, they would intercept most unauthorized transactions while rarely disrupting their innocent customers. 

This project reinforced practical lessons about dealing with imbalanced datasets and the realities of machine learning in production. I learned that deep learning isn't always the optimal solution for tabular data; sometimes a structured XGBoost model is exactly what is needed. Most importantly, my novel contribution proved that thoughtful feature engineering is often more highly correlated with performance improvements than simply selecting more complex model architectures. 

\section{Novel Contributions}

When brainstorming for my novel contribution, I realized that analyzing a single transaction independently limits contextual understanding. For instance, purchasing a \$500 television might appear statistically normal. However, purchasing a \$500 television at 2 AM, when it is the user's 15th transaction within the past hour, presents a highly suspicious pattern. 

My original idea was to engineer \textbf{Leakage-Safe Rolling Temporal Behavioral Features}. Instead of just passing the raw dataset variables, I authored a script in `data\_utils.py` that reconstructs the chronological history of every single credit card (`cc\_num`) to mathematically simulate account "velocity."

I engineered four new specific features:
\begin{enumerate}
    \item \textbf{time\_since\_last\_tx}: By sorting the data temporally and grouping by credit card number, I calculated the exact elapsed seconds since the user's previous transaction. Fraudsters conventionally swipe stolen credentials rapidly before the institution triggers a freeze.
    \item \textbf{rolling\_tx\_count\_1h}: Utilizing a pandas 1-hour rolling window, I accumulated the total transaction count for the preceding 60 minutes.
    \item \textbf{rolling\_tx\_count\_24h}: Similarly, a 24-hour aggregate transaction count.
    \item \textbf{rolling\_amt\_mean\_24h}: I calculated the rolling average expenditure for that card over the previous 24 hours. A transaction representing a massive standard deviation spike above this rolling mean acts as an anomaly flag.
\end{enumerate}

These engineered features were objectively validated as highly impactful. Looking at our summary table, adding these four variables increased XGBoost's F1 score from \textbf{0.674 to 0.892} ($\sim$32\% improvement), Random Forest's F1 from \textbf{0.466 to 0.723} ($\sim$55\% improvement), and the Neural Network's F1 from \textbf{0.759 to 0.874}. This quantitatively proved that providing models with moving averages over an entity's behavior is functionally superior to simple point-in-time attributes. 

Implementing this was technically demanding due to the risk of \textbf{data leakage}. If I included the current transaction's parameters within the rolling window calculations, the model would unfairly access future information during training. I utilized the `closed="left"` parameter within the pandas rolling window function to restrict the aggregation, ensuring the model only references the strict historical activity leading up to the exact second of the current transaction. This approach verified that my mathematical behavior profiling genuinely improved the baseline models without analytical cheating.

\section{Team Member Contribution}
As this was an individual project, I (Nathaniel Wetzel) was solely responsible for all components over the entire machine learning pipeline. This included problem definition and literature review, executing exploratory data analysis, designing a leakage-free feature engineering algorithm, orchestrating the Python training loop (with normalization pipelines), generating the statistical plots, and authoring the ultimate project report. All 4 algorithms and quantitative analyses were constructed and documented personally. 

\end{document}
